\chapter{Numbers and Operators}\label{ch-numbers}

Zeno has two kinds of plain number. An \verb|int| is a whole number, of
64 bits, from $-9223372036854775808$ to $9223372036854775807$. A
\verb|float| is a number with a fraction, the 64-bit floating point
number that every processor with a floating point unit has. Chapters
later in the book add numbers with a unit and whole numbers with a
smaller size. This chapter is the plain two, and the operators that
work on them.

\section{Writing numbers}

\begin{zeno}
print(42, -7, 1_000_000)          # _ can split the digits, anywhere
print(0xFF, 0b1010_0101)          # hexadecimal and binary
print(3.14, -0.5, 2.5e-7, 1e10)   # floats
\end{zeno}

A float always prints with a decimal point or an exponent, so that it
can never be taken for an int: \verb|3.0| prints as \verb|3.0|, not
\verb|3|.

\section{Arithmetic}

\begin{center}
\begin{tabular}{lll}
\toprule
Operator & Means & Example \\
\midrule
\verb|+| & add & \verb|7 + 2| is \verb|9| \\
\verb|-| & subtract, or negate & \verb|7 - 2| is \verb|5| \\
\verb|*| & multiply & \verb|7 * 2| is \verb|14| \\
\verb|/| & divide, always giving a float & \verb|7 / 2| is \verb|3.5| \\
\verb|//| & divide and round down & \verb|7 // 2| is \verb|3| \\
\verb|%| & the remainder of \verb|//| & \verb|7 % 2| is \verb|1| \\
\verb|**| & power & \verb|7 ** 2| is \verb|49| \\
\bottomrule
\end{tabular}
\end{center}

\begin{zeno}
print(7 + 2, 7 - 2, 7 * 2, 7 / 2, 7 // 2, 7 % 2, 7 ** 2)
print(10 / 5, 2 ** 0.5)
\end{zeno}

\subsection{Division}

\verb|/| always gives a float, even when the numbers divide exactly.
To divide whole numbers and keep a whole number, say so with
\verb|//|.

\begin{why}
In C, \verb|7 / 2| is \verb|3| and \verb|7.0 / 2| is \verb|3.5|: one
operator does two different things depending on what it is given, and
the remainder vanishes without a word. In Zeno the operator decides
what happens, not the values. (Decision 2.1.)
\end{why}

\verb|//| rounds \emph{down}, toward minus infinity, and \verb|%| takes
the sign of the number on its right. Together they always satisfy
\verb|(a // b) * b + a % b == a|:

\begin{zeno}
print(-7 // 2, -7 % 2, 7 % -2)
print(5.5 // 2, 5.5 % 2)
\end{zeno}

That is what makes \verb|i % n| a safe position in a list of length
\verb|n| for any \verb|i|, negative or not.

Dividing by zero is an error, with every division operator and with
both kinds of number:

\begin{zeno}[fails=yes]
print(1 / 0)
\end{zeno}

\subsection{Powers}

\verb|2 ** 10| is an int. A whole number to a negative power would be
a fraction, so it is refused, and the message says what to write:

\begin{zeno}[fails=yes]
print(2 ** -1)
\end{zeno}

\begin{why}
If \verb|2 ** -1| were \verb|0.5|, the type of \verb|a ** b| would
depend on the \emph{value} of \verb|b|, which nothing can know before
the program runs. The operator decides the type, not the values.
(Decision 2.3a.)
\end{why}

\subsection{Overflow}

An \verb|int| that would go past its largest or smallest value is an
\verb|OverflowError|, never a wrong number:

\begin{zeno}[fails=yes]
let big = 9223372036854775807
print(big + 1)
\end{zeno}

\begin{why}
In C, an overflow either wraps round without a word or is undefined.
A wrong number that looks plausible is the worst kind of mistake.
When a program really does want a number to wrap round, as a checksum
does, it says so with \verb|wrap|, which
\hyperref[ch-sized-numbers]{Sized Numbers and Bits} explains.
(Decision 2.2.)
\end{why}

A float that grows too large becomes \verb|inf|, as floats do
everywhere:

\begin{zeno}
print(1e308 * 10, -1e308 * 10)
\end{zeno}

\section{Mixing ints and floats}

An int and a float can meet in arithmetic, and the answer is a float:

\begin{zeno}
print(1 + 0.5, 3 * 1.5, 1 == 1.0)
\end{zeno}

Nothing else changes kind by itself. Text is never turned into a
number, or a number into text, without being asked:

\begin{zeno}[fails=yes]
print("3" + 4)
\end{zeno}

\begin{why}
Making an int a float in arithmetic is what people expect from
arithmetic. Every other automatic change is a guess at what the
programmer meant, and JavaScript shows where those guesses lead:
\verb|"3" + 4| is \verb|"34"| there. (Decision 2.3.)
\end{why}

\section{Changing kind on purpose}

\begin{center}
\begin{tabular}{lp{8cm}}
\toprule
Written & Gives \\
\midrule
\verb|int(x)| & a whole number: from text, or from a float by dropping the fraction \\
\verb|float(x)| & a float: from text or from an int \\
\verb|str(x)| & the value as text, as \verb|print| shows it \\
\bottomrule
\end{tabular}
\end{center}

\begin{zeno}
print(int("42"), int(" 7 "), int(9.99), int(-9.99))
print(float("1e3"), float(3), str(2.50) + "!")
\end{zeno}

Text that is not a number is a \verb|ValueError|, which a program can
catch (\hyperref[ch-errors]{Errors}):

\begin{zeno}[fails=yes]
print(int("forty"))
\end{zeno}

\section{Comparing}

\begin{center}
\begin{tabular}{ll}
\toprule
Operator & Means \\
\midrule
\verb|==|, \verb|!=| & equal, not equal \\
\verb|<|, \verb|<=| & less than, less than or equal \\
\verb|>|, \verb|>=| & greater than, greater than or equal \\
\bottomrule
\end{tabular}
\end{center}

A comparison gives a \verb|bool|. Comparisons \emph{chain}, as they do
in mathematics: \verb|0 <= i < n| means \verb|0 <= i and i < n|, with
\verb|i| worked out once.

\begin{zeno}
let i = 3
print(0 <= i < 10, 1 < 2 < 3, 3 > 2 > 2)
\end{zeno}

\begin{why}
In C, \verb|0 <= i < n| means \verb|(0 <= i) < n|, which compares true
or false with a number, compiles without a word, and is a classic
mistake. Refusing chains altogether, as Rust does, is safe but loses
the readable form. Zeno gives the line the meaning it looks like it
has. (Decision 1.16.)
\end{why}

Comparing two things that can never be equal is a mistake, and Zeno
says so instead of quietly answering \verb|false|. It is nearly always
a number that was read as text and never turned into one:

\begin{zeno}[fails=yes]
let typed = "1"
print(typed == 1)
\end{zeno}

Floats are kept in binary, so most decimal fractions are close to
what you wrote but not exactly it. That is true in every language, and
it is why floats are compared with care:

\begin{zeno}
print(0.1 + 0.2, 0.1 + 0.2 == 0.3)
print(abs((0.1 + 0.2) - 0.3) < 1e-9)
\end{zeno}

\section{True and false}

A \verb|bool| is \verb|true| or \verb|false|, and three words work on
them: \verb|and|, \verb|or| and \verb|not|.

\begin{zeno}
let warm = true
let raining = false
print(warm and not raining, warm or raining, not warm)
\end{zeno}

\verb|and| and \verb|or| stop as soon as the answer is known: in
\verb|a and b|, \verb|b| is not worked out at all when \verb|a| is
false. That makes \verb|len(xs) > 0 and xs[0] == 1| safe for an empty
list.

A condition must be a real \verb|bool|. A number, a list or a piece of
text is not true or false by itself:

\begin{zeno}[fails=yes]
let items = [1, 2]
if items {
    print("not empty")
}
\end{zeno}

\begin{why}
"Truthiness" means remembering which values count as false: is it
\verb|0|, \verb|""|, \verb|[]|, \verb|none|, \verb|0.0|? The answer
differs between languages and causes mistakes such as treating a real
reading of \verb|0| as no reading at all. Saying what you test,
\verb|len(items) > 0|, is clearer. (Decision 2.4.)
\end{why}

\section{Changing a variable: the = forms}

Every operator that takes two numbers has a form that changes a
variable, \verb|x op= y|, which means exactly \verb|x = x op y|:

\begin{center}
\verb|+=  -=  *=  /=  //=  %=  **=  &=  |\verb!|=!\verb|  ^=  <<=  >>=|
\end{center}

\begin{zeno}
var n = 10
n -= 3        # 7
n *= 2        # 14
n //= 3       # 4
n **= 2       # 16
print(n)
var f = 1.0
f /= 4
print(f)
\end{zeno}

The variable is read first and the right-hand side worked out second,
as in \verb|x = x + y|. \verb|/=| on an int is refused, because
\verb|/| gives a float and an int variable cannot hold one: use
\verb|//=|.

\begin{zeno}[fails=yes]
var half = 9
half /= 2
\end{zeno}

\section{Which operator goes first}

When operators are mixed, some bind tighter than others, as in school
arithmetic: \verb|2 + 3 * 4| is \verb|14|. From the loosest to the
tightest:

\begin{center}
\begin{tabular}{ll}
\toprule
Operators & \\
\midrule
\verb|or| & loosest \\
\verb|and| & \\
\verb|not| & \\
\verb|==  !=  <  <=  >  >=  in  is| & comparisons, which chain \\
\verb!|! & bit or \\
\verb|^| & bit exclusive or \\
\verb|&| & bit and \\
\verb|<<  >>| & shifts \\
\verb|+  -| & \\
\verb|*  /  //  %| & \\
\verb|-x  ~x| & negation, bit not \\
\verb|**| & tightest \\
\bottomrule
\end{tabular}
\end{center}

\begin{zeno}
print(2 + 3 * 4 ** 2, (2 + 3) * 4, -2 ** 2)
\end{zeno}

Brackets always win, and they are the best way to make a long
expression plain to the reader. The bit operators are in
\hyperref[ch-sized-numbers]{Sized Numbers and Bits}, and the whole
table again, with every operator, is \hyperref[app-operators]{Operators}.

\section{Built-in functions for numbers}

\begin{zeno}
print(abs(-4.5), min(3, 1, 2), max(2, 8.5), sum([1, 2, 3]))
print(min([5, 3, 9]), max([5, 3, 9]))
\end{zeno}

\verb|min| and \verb|max| take either several numbers or one list.
Square roots, rounding, logarithms and the trigonometric functions are
in the \verb|math| module (\hyperref[ch-math]{The math Module}).

\begin{summary}
\verb|int| is a 64-bit whole number, \verb|float| a 64-bit floating
point number. \verb|/| always gives a float and \verb|//| rounds down.
An int that overflows is an error, never a wrong number. Ints and
floats mix in arithmetic and nowhere else. Comparisons chain,
conditions must be real \verb|bool|s, and every operator on two
numbers has an \verb|=| form.
\end{summary}

\begin{tryit}
\begin{enumerate}
\item Work out \verb|-17 // 5| and \verb|-17 % 5| on paper, then check
  with Zeno. Check that \verb|(a // b) * b + a % b == a| holds.
\item What is the largest power of two that fits in an \verb|int|?
  Find it by trying \verb|2 ** n| for larger and larger \verb|n|.
\item Write \verb|if 1 <= x <= 10 and x % 2 == 0| for a few values of
  \verb|x|. Put brackets where Zeno would put them.
\item Print \verb|1 / 3 * 3|. Is it exactly 1? Now try
  \verb|0.1 * 3|.
\end{enumerate}
\end{tryit}
